\documentclass[10pt,a4paper]{amsart}
\usepackage[margin=19mm]{geometry}
\usepackage{amsmath,amssymb,mathtools}
\usepackage[scr=rsfs,cal=euler]{mathalfa}
\usepackage{xfrac}
\usepackage[unicode,hidelinks,bookmarks=false]{hyperref}
\newcommand{\ie}{i.e.\ }
\newcommand{\cf}{cf.\ }
\newcommand{\resp}{respectively\ }
\newcommand{\Subsection}{\subsection}
\providecommand{\nobreakdash}{\nobreak\hbox{-}}
\setlength{\emergencystretch}{2em}
\multlinegap0pt
\usepackage{amssymb}
\newtheorem{lemma}{Lemma}
\newtheorem{proposition}{Proposition}
\newcommand{\T}{\mathbb{T}}
\title{A Resolution of the de Bruijn--Erd\H{o}s Consecutive-Gap Problem}
\author{Samuel Korsky}
\date{Source: arXiv:2609.07196v2 (CC BY 4.0)}
\begin{document}
\maketitle

\noindent\emph{Source and licence.} A Resolution of the de Bruijn--Erd\H{o}s Consecutive-Gap Problem. Version arXiv:2609.07196v2 (CC BY 4.0), distributed by arXiv under the Creative Commons Attribution 4.0 International licence, http://creativecommons.org/licenses/by/4.0/, as recorded in the arXiv licence metadata for this exact version and shown on the abstract page https://arxiv.org/abs/2609.07196v2. Full licence notice: \texttt{licenses/arxiv-2609.07196-CC-BY-4.0.txt}.\par\smallskip

\noindent\textit{Editorial scope.} Counted unit: the article's proposition prop:short (source0.tex lines 378--395). The article's complete original proof of it is delivered with it (source0.tex lines 396--490, one \texttt{proof} environment of 95 lines). No mathematics is written, repaired or re-derived: the statement and the proof are the article's own lines, in the article's own notation, and no retained line is substituted or re-flowed.  Retained uncounted as the source-local context the proof needs: lines 239--242; lines 260--275.  Nothing was omitted from any retained span, so the package carries no omission record.  No retained line contains a citation, so the wrapper carries no bibliography and no bibliography entry.  No retained line contains a figure environment, an image-inclusion command or drawing code, so no image is delivered and the package declares no asset dependency.  Editorial formatting only, in addition: an AMS \texttt{amsart} preamble that restates the article's own declarations and macros used by the retained lines, namely the environments \texttt{lemma}, \texttt{proposition} on separate counters, together with the standard packages those lines need.  The retained environments print in this document as Lemma 1, Proposition 1 in order of appearance, whereas the article prints the same items with the numbering of its own full document.  The complete native source of the article as distributed in the arXiv e-print for this exact version is bundled unchanged as \texttt{sources/209609/source0.tex}.
\begin{equation}\label{eq:span-assumption}
  \frac{r-a_t}{t}\leq S_i(t)\leq\frac{r+b_t}{t},
  \qquad a_t+b_t\leq A,
\end{equation}

\begin{lemma}\label{lem:comparison}
Fix $D,E>0$ and an integer $k\geq1$, and put $q=E/(kr)$. If $q<1$, then,
for all sufficiently large $t$,
\begin{align}
  U_t(D)
  &\leq
  \left(1+\frac{3kA}{D}\right)(1+q)\,
  U_{(1+q)t}(E),\label{eq:comparison-upper}\\
  V_t(D)
  &\geq
  \left(1-q-\frac{kA}{D}(3-q)\right)_+
  V_{(1-q)t}(E).\label{eq:comparison-lower}
\end{align}
For fixed $E$ and $k$, the time threshold can be chosen uniformly for $D$
in any bounded range.
\end{lemma}

\begin{proposition}\label{prop:short}
There are constants $C_0,C_1>0$ with the following property. Suppose that
\eqref{eq:span-assumption} holds for all sufficiently large $t$, with
$A\geq1$ and $r\geq C_0A$. Set
\[
  \Lambda=\log(r/A),
  \qquad
  S=\frac{\sqrt{Ar}}{\Lambda^2}.
\]
Then, for every sufficiently large $t$,
\begin{equation}\label{eq:short-count}
  \sup_{x\in\T}\sup_{0\leq D\leq S}
  \left|N_t((x,x+D/t])-D\right|
  \leq3A+\frac{C_1A}{\Lambda}.
\end{equation}
The time threshold may depend on $r,A$, and the point sequence, but not
on $x$ or $D$.
\end{proposition}

\begin{proof}
Put
\[
  \theta=\sqrt{A/r},
  \qquad
  K=\sqrt{Ar}.
\]
We take $C_0$ sufficiently large throughout the proof; in particular,
$\theta\leq1/12$ and $K<r-A$.
\paragraph{The terminal bounds.}
An interval of length $(r-A)/t$ contains at most $r$ points. Indeed,
$r+1$ points in such a half-open interval would give an $r$-span of
strictly smaller length, contrary to~\eqref{eq:span-assumption}. Similarly,
an interval of length $(r+A)/t$ contains at least $r$ points: start at the
last point at or before its left endpoint and use the upper span bound.
Therefore
\begin{equation}\label{eq:terminal}
  U_t(r-A)\leq\frac r{r-A},
  \qquad
  V_t(r+A)\geq\frac r{r+A}
\end{equation}
for all sufficiently large $t$.
\paragraph{Transferring the terminal bounds to $K$.}
Construct a list of scales by starting at $K$ and doubling until reaching
$r-A$ for the upper estimate, and $r+A$ for the lower estimate. Shorten
the last step in each list to end at its specified terminal scale.
The estimates propagate from these known terminal bounds back down to
$K$. For two adjacent scales $D,E$, we have
\[
  K\leq D\leq E\leq2D.
\]
Choose $k=\lceil D/K\rceil$ in Lemma~\ref{lem:comparison}. Since
$k\leq2D/K$,
\[
  q=\frac{E}{kr}\leq2\theta,
  \qquad
  \frac{3kA}{D}\leq6\theta.
\]
The upper multiplier is at most
\[
  (1+6\theta)(1+2\theta)\leq1+9\theta,
\]
and the lower multiplier is at least $1-8\theta>0$.
Let $h_U$ and $h_V$ be the numbers of comparisons in the two lists. Both
are at most $C\Lambda$ for an absolute constant $C$. The terminal times
are obtained from $t$ by finite products of positive factors $1+q$ or
$1-q$. Thus, for sufficiently large $t$, every comparison and both
terminal bounds apply. Iteration gives the explicit estimates
\[
  U_t(K)\leq\frac r{r-A}(1+9\theta)^{h_U},
  \qquad
  V_t(K)\geq\frac r{r+A}(1-8\theta)^{h_V}.
\]
Since $A/r=\theta^2$ and $\theta\Lambda\to0$ as $r/A\to\infty$, these
imply
\begin{equation}\label{eq:middle}
  U_t(K)\leq1+O(\theta\Lambda),
  \qquad
  V_t(K)\geq1-O(\theta\Lambda)
\end{equation}
for all sufficiently large $t$. The implied constants are independent
of $r,A$, and the point sequence.
\paragraph{Transferring the bounds at $K$ to shorter intervals.}
Apply Lemma~\ref{lem:comparison} once more, with $E=K$ and $k=1$, so that
$q=\theta$. For $D>0$ and $I=(x,x+D/t]$, it gives
\begin{align*}
  N_t(I)&\leq
    (D+3A)(1+\theta)\,U_{(1+\theta)t}(K),\\
  N_t(I)&\geq
    \bigl(D(1-\theta)-(3-\theta)A\bigr)_+\,
    V_{(1-\theta)t}(K).
\end{align*}
Using~\eqref{eq:middle}, we obtain
\begin{equation}\label{eq:additive-error}
  |N_t(I)-D|
  \leq3A+O\bigl((D+A)\theta\Lambda\bigr).
\end{equation}
For the lower estimate, if the expression inside the positive part is
nonpositive, then $D\leq3A+O(A\theta)$ and the trivial bound $N_t(I)\geq0$
suffices. Otherwise the lower bound in~\eqref{eq:middle} applies directly.
The scale $S=K/\Lambda^2$ is chosen to make the remaining error
$O(A/\Lambda)$. Indeed, for $D\leq S$,
\[
  D\theta\Lambda\leq\frac A\Lambda,
  \qquad
  A\theta\Lambda=O(A/\Lambda),
\]
where the second estimate follows from $\theta\Lambda^2\to0$.
This proves~\eqref{eq:short-count} after fixing sufficiently large absolute
$C_0$ and $C_1$.
The scale lists use only finitely many predetermined comparison times,
and the last comparison uses $(1\pm\theta)t$, independently of $D$.
Consequently a single time threshold works for every $x$ and
$0<D\leq S$. The case $D=0$ is immediate.
\end{proof}

\end{document}
